% Part: first-order-logic
% Chapter: lindstrom
% Section: abstract-logics

\documentclass[../../../include/open-logic-section]{subfiles}

\begin{document}

\olfileid{mod}{lin}{alg}
\olsection{నైరూప్య తర్కాలు}

\begin{defn}
\emph{నైరూప్య తర్కం} అంటే $\tuple{L, \models_L}$ అనే జత; ఇక్కడ $L$
అనేది ప్రతి \tetoken{భాషకు}{!!{language}}~$\Lang{L}$,
\tetoken{వాక్యాల}{!!{sentence}s} సమితి $L(\Lang{L})$ను కేటాయించే
ప్రమేయం; $\models_L$ అనేది $\Lang{L}$కు చెందిన
\tetoken{నిర్మాణాలకు}{!!{structure}s}, ఆ
\tetoken{భాష}{!!{language}} యొక్క $L(\Lang{L})$లోని
\tetoken{మూలకాలకు}{!!{element}s} మధ్య సంబంధం. ముఖ్యంగా,
$\tuple{F, \models}$ అనేది మామూలు మొదటిస్థాయి తర్కం; అంటే, $F$ అనేది
\tetoken{భాషకు}{!!{language}}~$\Lang{L}$, $\Lang{L}$లోని సంకేతాలను
ఉపయోగించి నిర్మించిన మొదటిస్థాయి \tetoken{వాక్యాల}{!!{sentence}s}
సమితిని కేటాయించే ప్రమేయం; $\models$ అనేది మొదటిస్థాయి తర్కపు
సంతృప్తి సంబంధం.
\sourcecorrection{OLTEMODLIN-016}{మొదటిస్థాయి వాక్యాలు
$\Lang{L}$లోని “స్థిరసంకేతాల” నుంచి నిర్మితమవుతాయని ఆంగ్ల మూలం
అంటుంది. కానీ ఈ అధ్యాయపు సంబంధాత్మక భాషల్లో విధేయ సంకేతాలు,
వ్యక్తిగత స్థిరసంకేతాలు రెండూ ఉంటాయి; పరమాణు సూత్రాలకు విధేయ
సంకేతాలు అవసరం. అందువల్ల ఉద్దేశించిన పూర్తి “సంకేతాలు” అనే
పరిధిని తెలుగులో స్పష్టం చేశాం.}
\end{defn}

ఇక్కడ కూడా, ఇచ్చిన \tetoken{భాషకు}{!!{language}} మొదటిస్థాయి
తర్కంలో వాడిన \tetoken{నిర్మాణం}{!!{structure}} అనే భావననే వాడుతున్నాం.
అయితే \tetoken{వాక్యాలు}{!!{sentence}s} $\Lang{L}$లోని మౌలిక
సంకేతాల నుంచి మామూలు విధంగా నిర్మితమవుతాయని గానీ, $\models_L$
సంబంధం మొదటిస్థాయి తర్కంలోలాగే పునరావృత్తంగా నిర్వచితమవుతుందని గానీ
ముందే ఊహించం. ఉదాహరణకు, పూర్తిగా సాధారణమైన ఈ నిర్వచనం
$\tuple{L,\models_L}$లోని \tetoken{వాక్యాలు}{!!{sentence}s} అనంత
పొడవైన సంయోగాలు లేదా వియోగాలను, $\lexists$, $\lforall$ కాక వేరే
పరిమాణకాలను (ఉదాహరణకు, “\dots~అయ్యే $x$లు అనంత సంఖ్యలో ఉన్నాయి”),
లేదా బహుశా అనంత పొడవైన పరిమాణక పూర్వపదాలను కలిగిన సందర్భాలను కూడా
ఇముడ్చాలనే ఉద్దేశంతో ఉంది. $L(\Lang{L})$లోని
“\tetoken{వాక్యాలు}{!!{sentence}s}” మామూలు మొదటిస్థాయి తర్కపు
\tetoken{వాక్యాలు}{!!{sentence}s} కావలసిన అవసరం లేదని నొక్కి చెప్పడానికి,
ఈ అధ్యాయంలో వాటిపై విస్తరించడానికి \tetoken{చరాలు}{!!{variable}s}
$!E$, $!F$,~\dots వాడి, మామూలు మొదటిస్థాయి
\tetoken{సూత్రాలకు}{!!{formula}s} $!A$, $!B$,~\dotsను కేటాయిస్తాం.

\begin{defn}
$\Mod(L){!E}$ అనేది $\Setabs{\Struct{M}}{\Struct{M}
  \models_L !E}$ అనే తరగతిని సూచించనివ్వండి. \tetoken{భాషను}{!!{language}}
స్పష్టంగా చూపవలసి వస్తే $\Mod[L](L){!E}$ అని రాస్తాం.
$\Lang{L}$కు చెందిన రెండు \tetoken{నిర్మాణాలు}{!!{structure}s}
$\Struct{M}$, $\Struct{N}$లలో $L(\Lang{L})$కు చెందిన అవే
\tetoken{వాక్యాలు}{!!{sentence}s} సత్యమైతే, అవి $\tuple{L,
\models_L}$లో \emph{వాక్య తుల్యాలు} అంటాం; దీనిని
$\Struct{M} \elemequiv[L] \Struct{N}$గా రాస్తాం.
\end{defn}

\begin{defn}
\tetoken{భాష}{!!{language}}~$\Lang{L}$కు చెందిన నైరూప్య తర్కం
$\tuple{L,\models_L}$ కింది ధర్మాలను నెరవేరిస్తే దాన్ని
\emph{సాధారణం} అంటాం:
\begin{enumerate}
\item (\emph{$L$-ఏకదిశత}) \tetoken{భాషలు}{!!{language}s}
  $\Lang{L}$, $\Lang{L'}$లకు $\Lang{L} \subseteq \Lang{L'}$ అయితే,
  $L(\Lang{L}) \subseteq L(\Lang{L'})$.
\item (\emph{విస్తరణ ధర్మం}) ప్రతి $!E \in L(\Lang{L})$కు
  $\Lang{L}$లో ఒక \emph{పరిమిత} ఉపసమితి $\Lang{L'}$ ఉండి,
  $!E \in L(\Lang{L'})$ కూడా అవుతుంది; అంతేకాక $\Struct{M}
  \models_L !E$ సంబంధం, $\Struct{M}$ను $\Lang{L'}$కు సంకుచితం చేసిన
  రూపంపై మాత్రమే ఆధారపడుతుంది. అంటే, $\Struct{M}$, $\Struct{N}$లను
  $\Lang{L'}$కు సంకుచితం చేసిన రూపాలు ఒకటే అయితే, $\Struct{M}
  \models_L !E$ అవుతుంది అప్పుడే $\Struct{N} \models_L !E$.
  \sourcecorrection{OLTEMODLIN-004}{ఆంగ్ల నిర్వచనం $!E$ సత్యత ఒక
  పరిమిత ఉపభాషకు సంకుచితం చేసిన రూపంపై ఆధారపడుతుందని మాత్రమే
  చెబుతుంది; తరువాతి రెండు నిరూపణలు అదే ధర్మం నుంచి
  $!E\in L(\Lang{L'})$ను కూడా ఉపయోగిస్తాయి. మొదటిస్థాయి వాక్యానికి
  ఉండే పరిమిత పదజాల ధర్మానికి, తరువాతి ప్రయోగాలకు అనుగుణంగా ఆ
  తప్పిపోయిన సభ్యత్వ నిబంధనను చేర్చాం.}
\item (\emph{సమరూపతా ధర్మం}) $\Struct{M} \models_L !E$,
  $\Struct{M} \simeq \Struct{N}$ అయితే, $\Struct{N} \models_L !E$
  కూడా.
\item (\emph{పునర్నామకరణ ధర్మం}) పునర్నామకరణ కింద $\models_L$
  సంబంధం సంరక్షితమవుతుంది: \tetoken{భాష}{!!{language}}~$\Lang{L}$లోని
  ప్రతి సంకేతం $P$ని అదే అరిటీ గల $P'$తో, ప్రతి స్థిరసంకేతం $c$ను
  వేరొక స్థిరసంకేతం $c'$తో మార్చి $\Lang{L}'$ను పొందితే, ప్రతి
  \tetoken{నిర్మాణం}{!!{structure}}~$\Struct{M}$కు, ప్రతి
  \tetoken{వాక్యం}{!!{sentence}}~$!E$కు, $\Struct{M} \models_L !E$
  అవుతుంది అప్పుడే $\Struct{M}' \models_L !E'$. ఇక్కడ
  $\Struct{M}'$ అనేది $\Struct{M}$కు అనురూపమైన
  $\Lang{L}'$-\tetoken{నిర్మాణం}{!!{structure}}, అలాగే $!E' \in
  L(\Lang{L}')$.
  \sourcecorrection{OLTEMODLIN-001}{పునర్నామకృత
  $\Struct{M}'$ను ఆంగ్ల మూలం పొరపాటుగా $\Lang{L}$కు అనురూపమైన
  నిర్మాణం అంటుంది. ప్రదర్శిత ద్విసంబంధం కోరినట్లుగా అది అసలు
  నిర్మాణం $\Struct{M}$కు అనురూపమైన $\Lang{L}'$-నిర్మాణమని
  సరిచేశాం.}
\item (\emph{బూలియన్ ధర్మం}) నైరూప్య తర్కం $\tuple{L,\models_L}$
  బూలియన్ సంయోజకాల కింద కింది అర్థంలో సంవృతం: ప్రతి $!E \in
  L(\Lang{L})$కు, $\Struct{M} \models_L !F$ అవుతుంది అప్పుడే
  $\Struct{M} \not\models_L !E$ అయ్యే $!F \in L(\Lang{L})$ ఉంటుంది;
  ప్రతి $!E$, $!F$లకు $\Mod(L){!G} = \Mod(L){!E} \cap
  \Mod(L){!F}$ అయ్యే $!G$ ఉంటుంది. ఇంకా, స్థిరసంకేతాలతో ఏర్పడిన
  మొదటిస్థాయి పరమాణు \tetoken{సూత్రాలను}{!!{formula}s} వాక్యాలుగా
  తీసుకున్నప్పుడు అవి ఈ తర్కంలో ఉండి, వాటికి సాధారణ మొదటిస్థాయి సంతృప్తి
  సంబంధమే వర్తిస్తుంది; మిగతా బూలియన్ సంయోజకాలకూ పై విధమైన సంవృతత
  ఉంటుంది.
  \sourcecorrection{OLTEMODLIN-018}{నైరూప్య తర్కంలో వాక్యాలను,
  నిర్మాణాలతో వాటి సంతృప్తి సంబంధాన్ని మాత్రమే నిర్వచించిన తరువాత
  ఆంగ్ల మూలం “పరమాణు సూత్రాల విషయంలోనూ ఇదే” అంటుంది; స్వేచ్ఛా
  చరాలున్న సూత్రాలకు ఇక్కడ అర్థవిచారం నిర్వచితం కాలేదు. తరువాత
  మొదటిస్థాయి తర్కాన్ని ఆగమనంతో ఇమిడ్చడానికి అవసరమైన, స్థిరసంకేతాలతో
  ఏర్పడిన పరమాణు వాక్యాల సాధారణ అర్థవిచారాన్ని స్పష్టంగా తెలిపాం.}
\item (\emph{పరిమాణక ధర్మం}) $\Lang{L}$లోని ప్రతి స్థిరసంకేతం $c$కు,
  ప్రతి $!E \in L(\Lang{L})$కు కింది విధంగా ఉండే $!F \in
  L(\Lang{L'})$ ఉంటుంది:
  \[
  \Mod[L'](L){!F} = \Setabs{\Struct{M}}{\Expan{M}{a} \in
  \Mod[L](L){!E} \text{ అయ్యే ఏదో } a \in \Domain{M}},
  \]
  ఇక్కడ $\Lang{L'} = \Lang{L} \setminus \{c\}$; $\Expan{M}{a}$ అనేది
  $c$కు $a$ను కేటాయిస్తూ $\Struct{M}$ను $\Lang{L}$కు చేసిన విస్తరణ.
  \sourcecorrection{OLTEMODLIN-019}{స్థిరసంకేతం $c$ను బంధించి
  భాష నుంచి తొలగించిన ఫలిత వాక్యం $!F$ను ఆంగ్ల మూలం
  $L(\Lang{L})$లోనే ఉంచుతుంది; కానీ వెంటనే దాని నమూనాల తరగతిని
  $\Lang{L'}=\Lang{L}\setminus\{c\}$కు సూచిస్తుంది. ఉద్దేశించిన
  $!F\in L(\Lang{L'})$ను పునరుద్ధరించాం.}
\item (\emph{సాపేక్షీకరణ ధర్మం}) \tetoken{ఒక వాక్యం}{!!a{sentence}}
  $!E \in L(\Lang{L})$, $\Lang{L}$లో లేని సంకేతాలు $R$, $c_1$,
  \dots, $c_n$ ఇవ్వబడ్డాయనుకుందాం. అప్పుడు
  \tetoken{ఒక వాక్యం}{!!a{sentence}}~$!F \in L(\Lang{L} \cup \{R,c_1,\ldots,c_n\})$గా ఉండి, దాన్ని $!E$కు
  $\Atom{R}{x, c_1, \dots c_n}$పై \emph{సాపేక్షీకరణ} అంటాం; ప్రతి
  \tetoken{నిర్మాణం}{!!{structure}}~$\Struct{M}$కు
  \[
  \Expan{M}{X, b_1, \ldots, b_n} \models_L !F
    \text{ అవుతుంది అప్పుడే } \Struct{N} \models_L !E,
  \]
  ఇక్కడ $\Struct{N}$ అనేది $\Struct{M}$ యొక్క, కింది
  \tetoken{వ్యక్తి క్షేత్రం}{!!{domain}} గల ఉపనిర్మాణం:
  $\Domain{N} = \Setabs{a\in \Domain{M}}{X(a, b_1, \dots, b_n)}$
  \sourcecorrection{OLTEMODLIN-002}{కొత్త సంకేతం $R$కు ఇంకా
  అర్థనిర్దేశం లేని అసలు $\Lang{L}$-నిర్మాణం $\Struct{M}$లోనే
  $\Assign{R}{M}$ను ఆంగ్ల మూలం వాడుతుంది. అదే వాక్యం విస్తరణలో
  $R$కు $X$ను అర్థనిర్దేశంగా ఇస్తుంది కాబట్టి, ఉద్దేశించిన
  $X(a,b_1,\ldots,b_n)$ తంతువును వాడాం.}
  చూపిన సమితి ఖాళీ కానిదై, $\Lang{L}$లోని అన్ని అసలు
  స్థిరసంకేతాల అర్థనిర్దేశాలను కలిగిన సందర్భాల్లోనే ఈ ఉపనిర్మాణాన్ని
  తీసుకుంటాం
  (\olref[bas][sub]{rem:substructure} చూడండి).
  \sourcecorrection{OLTEMODLIN-003}{ఈ గ్రంథంలో నిర్మాణపు వ్యక్తి
  క్షేత్రం ఖాళీ కానిది; ఈ అధ్యాయపు సంబంధాత్మక భాషల్లో వ్యక్తిగత
  స్థిరసంకేతాలు కూడా ఉండవచ్చు. యథేచ్ఛ $X$-తంతువు ఖాళీ కావచ్చు లేదా
  స్థిరసంకేతాల అర్థనిర్దేశాలను వదిలేయవచ్చు; అప్పుడు అది
  ఉపనిర్మాణం కాదు. అందువల్ల అవసరమైన ఖాళీకానితనం, స్థిరసంకేత-సంవృత
  నిబంధనలను స్పష్టం చేశాం.}
  $\Expan{M}{X, b_1, \ldots,b_n}$ అనేది $R$, $c_1$, \dots, $c_n$లకు
  వరుసగా $X$, $b_1,$ \dots, $b_n$లను అర్థనిర్దేశాలుగా ఇస్తూ
  $\Struct{M}$ను విస్తరించిన నిర్మాణం (ఇక్కడ $X \subseteq M^{n+1}$).
\end{enumerate}
\end{defn}

\begin{defn}
రెండు నైరూప్య తర్కాలు $\tuple{L_1, \models_{L_1}}$,
$\tuple{L_2, \models_{L_2}}$ ఇవ్వబడ్డాయనుకుందాం. ప్రతి
\tetoken{భాష}{!!{language}}~$\Lang{L}$కు, ప్రతి
\tetoken{వాక్యం}{!!{sentence}}~$!E \in L_1(\Lang{L})$కు
$\Mod[L](L_1){!E} = \Mod[L](L_2){!F}$ అయ్యే
\tetoken{ఒక వాక్యం}{!!a{sentence}}~$!F \in L_2(\Lang{L})$ ఉంటే,
రెండవ తర్కం మొదటిదానికన్నా \emph{కనీసం అంత వ్యక్తీకరణశక్తి గలది}
అంటాం; దీనిని $\tuple{L_1, \models_{L_1}} \leq
\tuple{L_2, \models_{L_2}}$గా రాస్తాం. $\tuple{L_1,
\models_{L_1}} \leq \tuple{L_2, \models_{L_2}}$, $\tuple{L_2,
\models_{L_2}} \leq \tuple{L_1, \models_{L_1}}$ రెండూ అయితే,
$\tuple{L_1, \models_{L_1}}$, $\tuple{L_2, \models_{L_2}}$లను
\emph{తుల్యమైనవి} అంటాం.
\end{defn}

\begin{rem}
మొదటిస్థాయి తర్కం, అంటే నైరూప్య తర్కం $\tuple{F, \models}$,
సాధారణమైనది. నిజానికి, పై ధర్మాల్లో చాలావరకు మొదటిస్థాయి తర్కానికి
నేరుగా వస్తాయి. విస్తరణ ధర్మం వాక్యపు పరిమిత పదజాలానికి, విస్తారతకు
దిగివస్తుందని, అలాగే
\tetoken{ఒక వాక్యం}{!!a{sentence}}~$!E$కు
$\Atom{R}{x, c_1, \dots, c_n}$పై సాపేక్షీకరణను, ప్రతి
\tetoken{ఉపసూత్రాన్ని}{!!{subformula}}~$\lforall[x][!F]$,
$\lforall[x][(\Atom{R}{x, c_1, \dots, c_n} \lif \ !F)]$తో మార్చి
పొందుతామని మాత్రమే గమనిద్దాం. అంతేకాక $\tuple{L,\models_L}$ సాధారణం
అయితే, మొదటిస్థాయి \tetoken{సూత్రాలపై}{!!{formula}s} ఆగమనం ద్వారా
  $\tuple{F, \models} \leq \tuple{L, \models_{L}}$ అని చూపవచ్చు. అందువల్ల,
సాధారణత్వాన్ని కోల్పోకుండా ప్రతి మొదటిస్థాయి
\tetoken{వాక్యం}{!!{sentence}} ప్రతి సాధారణ తర్కంలోనూ ఉంటుందని
అనుకోవచ్చు.
\end{rem}

\end{document}
